Se implementan dos arquitecturas: NCF, que modela cada par (usuario, película) por separado, y un Autocodificador, que reconstruye el perfil completo de un usuario en una sola pasada. Los ratings se normalizan a $[0,1]$ para ajustarlos a la activación sigmoide de salida, y se desnormalizan al calcular RMSE y MAE para comparar con los modelos anteriores.

\section{Neural Collaborative Filtering (NCF)}

El CF clásico predice el rating como el producto escalar de dos embeddings:

\begin{equation}
    \hat{r}_{ui} = \mathbf{p}_u^\top \mathbf{q}_i
\end{equation}

Esto es inherentemente lineal. NCF reemplaza ese producto por un MLP: los embeddings de usuario e ítem (32 dimensiones cada uno) se concatenan y pasan por capas densas $128 \to 64 \to 32$ con activación ReLU, hasta una salida sigmoide en $[0,1]$:

\begin{equation}
    \mathrm{ReLU}(x) = \max(0,\, x)
\end{equation}

La no-linealidad permite que la red aprenda interacciones entre embeddings que el producto escalar no puede representar.

\subsection{Regularización}

Cada capa densa va seguida de Dropout: 30\,\% en la primera capa y 20\,\% en la segunda. En entrenamiento, Dropout desactiva aleatoriamente una fracción de las neuronas en cada paso, forzando a la red a no depender de activaciones concretas. En inferencia, todas las neuronas están activas.

\subsection{Optimizador Adam}

El descenso de gradiente estándar actualiza todos los parámetros con el mismo ritmo fijo $\alpha$. El problema es que en una red neuronal los parámetros se comportan de forma muy distinta: algunos necesitan ajustes grandes porque apenas han cambiado, otros necesitan pasos pequeños porque llevan varias épocas oscilando. Un ritmo único es un compromiso malo para todos.

Adam lleva dos contadores por cada parámetro. El primero acumula la dirección media del gradiente:

\begin{equation}
    m_t = \beta_1 \cdot m_{t-1} + (1 - \beta_1) \cdot g_t
\end{equation}

Con $\beta_1 = 0.9$, el 90\,\% del valor viene del historial y solo el 10\,\% del gradiente del paso actual. Si el gradiente lleva varios pasos apuntando en la misma dirección, $m_t$ refleja eso con claridad. Si da un salto puntual por ruido, $m_t$ apenas se mueve.

El segundo acumula la variabilidad histórica del gradiente:

\begin{equation}
    v_t = \beta_2 \cdot v_{t-1} + (1 - \beta_2) \cdot g_t^2
\end{equation}

Elevar al cuadrado elimina el signo y amplifica las oscilaciones grandes. Con $\beta_2 = 0.999$, este contador tiene memoria larga. Un parámetro con gradientes grandes y erráticos acumula un $v_t$ alto, mientras que uno estable lo tiene bajo.

La actualización combina los dos:

\begin{equation}
    \theta_t = \theta_{t-1} - \frac{\alpha}{\sqrt{\hat{v}_t} + \varepsilon} \cdot \hat{m}_t
\end{equation}

El paso efectivo es $\alpha$ dividido entre $\sqrt{\hat{v}_t}$. Si un parámetro ha oscilado mucho, $v_t$ es grande, el denominador crece y el paso se achica. Si ha sido estable, $v_t$ es pequeño y el paso es mayor. El resultado es que cada parámetro tiene su propio ritmo de aprendizaje adaptado a su comportamiento. El $\varepsilon = 10^{-8}$ evita dividir entre cero. Los valores $\hat{m}_t$ y $\hat{v}_t$ son versiones corregidas para compensar el sesgo hacia cero que tienen ambas estimaciones al inicio del entrenamiento.

\section{Autocodificador colaborativo}

Cada usuario se representa como un vector de longitud igual al número de películas del catálogo: calificación normalizada en las posiciones que ha valorado, cero en el resto. El autocodificador aprende a comprimir ese vector en 64 dimensiones (codificador: $512 \to 64$) y reconstruirlo (decodificador: $64 \to 512 \to n_\text{películas}$). Las posiciones reconstruidas para ítems no vistos son las predicciones.

Con aproximadamente el 98\,\% de dispersión, calcular el MSE sobre todo el vector reconstruido haría que el gradiente lo dominasen los ceros, por lo que el modelo aprendería a predecir cero para todo. Por tanto, la pérdida se calcula solo sobre posiciones con calificación conocida:

\begin{equation}
    \mathcal{L} = \frac{\displaystyle\sum_{(u,i):\,r_{ui}\neq 0}(r_{ui} - \hat{r}_{ui})^2}
                       {\displaystyle\sum_{(u,i)} \mathbf{1}[r_{ui}\neq 0] + \varepsilon}
\end{equation}

donde $\varepsilon = 10^{-8}$ evita la división por cero.

El autocodificador genera las predicciones de todos los ítems de un usuario en una sola pasada. NCF necesita una pasada por par (usuario, ítem), lo que supone 10\,000 pasadas para un catálogo de 10\,000 películas. Los resultados sobre el conjunto de prueba son:

\begin{itemize}
    \item \textbf{NCF}: RMSE 0.9324, MAE 0.7031, tiempo 12.93\,s
    \item \textbf{Autocodificador}: RMSE 0.9553, MAE 0.7282, tiempo 12.88\,s
\end{itemize}

Los números son decepcionantes. Ambas redes son más lentas que KNN-User y peores en precisión que cualquiera de las dos variantes KNN. El Autocodificador es el modelo con peor RMSE de todos los implementados. La explicación no es que las arquitecturas sean malas (en datasets de mayor escala estas redes demuestran ser competitivas), sino que MovieLens Small tiene demasiado pocos datos para que los embeddings aprendan representaciones significativas. Con 610 usuarios y una media de 165 ratings por usuario, los embeddings de 32 dimensiones están sobreajustando al ruido del conjunto de entrenamiento en lugar de capturar preferencias reales.

Para usuarios con historial rico, ambos modelos producen recomendaciones distintas a las de KNN: emergen películas que no comparten géneros obvios con lo visto pero sí aparecen en patrones de co-valoración latentes.
